\section{进阶要求：观察 AI 编译器 MLIR 中端处理过程}

我们在 WSL2 Ubuntu 22.04 中使用 MLIR 14、\texttt{mlir-opt-14}、\texttt{mlir-translate-14} 和 Clang 14，先处理基础算术程序，再处理固定尺寸的矩阵乘法程序。我们依次查看高层 \texttt{linalg.matmul}、循环、控制流、LLVM Dialect、标准 LLVM IR 和 x86-64 目标文件。

\labfigure{figures/mlir/mlir_env.png}{MLIR 14 实验环境与工具版本}{fig:mlir-env}

\subsection{实验任务}

成员A负责基础 MLIR 程序和矩阵乘法前两步转换，执行 \texttt{canonicalize} 和 \texttt{convert-linalg-to-loops}；成员B负责核对 \texttt{convert-scf-to-std} 后的基本块和分支，以及 LLVM Dialect、标准 LLVM IR 与目标文件。我们共同整理各阶段的文件和终端输出。

\subsection{基础 MLIR 程序}

我们在 \texttt{mlir/test.mlir} 中定义两个 \texttt{i32} 常量和一次 \texttt{arith.addi}，再执行规范化转换，观察函数体的变化。

\lstinputlisting[
caption={基础 MLIR 程序 test.mlir},
label={lst:mlir-test}
]{../mlir/test.mlir}

先直接查看程序：

\begin{lstlisting}[language=bash]
mlir-opt-14 test.mlir
\end{lstlisting}

我们在输出中查看 \texttt{module}、\texttt{func}、\texttt{arith.constant} 和 \texttt{arith.addi}。图\ref{fig:mlir-basic}给出了该命令的输出。

\labfigure{figures/mlir/mlir_basic.png}{基础 MLIR 程序的输出}{fig:mlir-basic}

再执行 \texttt{canonicalize}：

\begin{lstlisting}[language=bash]
mlir-opt-14 test.mlir -canonicalize
\end{lstlisting}

执行规范化后，函数体保留 \texttt{return}。图\ref{fig:mlir-canonicalize}给出了转换前后的结果。

\labfigure{figures/mlir/mlir_canonicalize.png}{Canonicalization 前后的 MLIR 对比}{fig:mlir-canonicalize}

\subsection{Linalg Dialect 中的矩阵乘法}

我们在矩阵乘法程序 \texttt{mlir/matmul.mlir} 中定义 $2\times4$ 的 $A$、$4\times3$ 的 $B$ 和 $2\times3$ 的 $C$，并用一条 \texttt{linalg.matmul} 操作表示计算。

\lstinputlisting[
caption={Linalg Dialect 中的矩阵乘法程序 matmul.mlir},
label={lst:mlir-matmul}
]{../mlir/matmul.mlir}

\begin{lstlisting}[language=bash]
mlir-opt-14 matmul.mlir
\end{lstlisting}

我们在图\ref{fig:mlir-linalg-matmul}中查看矩阵乘法的输入、输出和维度。

\labfigure{figures/mlir/mlir_linalg_matmul.png}{Linalg Dialect 中的矩阵乘法表示}{fig:mlir-linalg-matmul}

\subsection{Linalg 到 SCF、MemRef 和 Arith 的 Lowering}

使用 \texttt{convert-linalg-to-loops} 后，我们将 \texttt{linalg.matmul} 展开为三层 \texttt{scf.for}。循环体使用 \texttt{memref.load} 读取 $A[i][k]$、$B[k][j]$ 和 $C[i][j]$，使用 \texttt{arith.mulf}、\texttt{arith.addf} 完成乘加，再由 \texttt{memref.store} 写回 $C[i][j]$。

\begin{lstlisting}[language=bash]
mlir-opt-14 matmul.mlir -convert-linalg-to-loops
mlir-opt-14 matmul.mlir -convert-linalg-to-loops -o matmul_loops.mlir
\end{lstlisting}

\labfigurecompact{figures/mlir/mlir_linalg_loops.png}{linalg.matmul Lowering 为三层 SCF 循环}{fig:mlir-linalg-loops}{0.62}

\labfigurecompact{figures/mlir/mlir_loops_saved.png}{matmul\_loops.mlir 的保存结果}{fig:mlir-loops-saved}{0.62}

我们在 \texttt{matmul\_loops.mlir} 中查看三层循环：最外层范围为 $0$ 到 $2$，中间层为 $0$ 到 $3$，最内层为 $0$ 到 $4$。原来的高层操作在此阶段展开为循环边界、显式内存访问和浮点乘加。

\subsection{SCF 到控制流形式的 Lowering}

MLIR 14 的 Pass 名称与资料中的写法存在差异。我们先通过帮助信息确认本机可用的转换 Pass，再使用 \texttt{convert-scf-to-std} 处理保存的循环程序。

\begin{lstlisting}[language=bash]
mlir-opt-14 --help | grep -E 'convert-(scf|memref|arith|std)'
mlir-opt-14 matmul_loops.mlir -convert-scf-to-std -o matmul_cfg.mlir
\end{lstlisting}

\labfigure{figures/mlir/mlir_pass_help.png}{MLIR 14 中可用的关键 Lowering Pass}{fig:mlir-pass-help}

我们在 \texttt{matmul\_cfg.mlir} 中查看三层循环转换出的基本块、\texttt{br}、\texttt{cond\_br} 和块参数；\texttt{arith.cmpi} 控制循环是否继续。\texttt{memref.load}、\texttt{arith.mulf}、\texttt{arith.addf} 和 \texttt{memref.store} 位于相应的基本块中。

\labfigure{figures/mlir/mlir_cfg.png}{SCF 结构化循环转为基本块和条件分支}{fig:mlir-cfg}

\subsection{转换为 LLVM Dialect}

在控制流转换完成后，我们继续将 Arith、MemRef 和 Standard 操作转换为 LLVM Dialect，并处理转换中的临时类型桥接：

\begin{lstlisting}[language=bash]
mlir-opt-14 matmul_cfg.mlir \
  -convert-arith-to-llvm -convert-memref-to-llvm \
  -convert-std-to-llvm -reconcile-unrealized-casts \
  -o matmul_llvm.mlir
\end{lstlisting}

我们在转换结果中查看 \texttt{llvm.func}、\texttt{llvm.br}、\texttt{llvm.cond\_br}、\texttt{llvm.getelementptr}、\texttt{llvm.load}、\texttt{llvm.fmul}、\texttt{llvm.fadd} 和 \texttt{llvm.store}。memref 访问在这一阶段写成指针、偏移和维度信息，矩阵元素地址由 \texttt{llvm.getelementptr} 计算。

\labfigure{figures/mlir/mlir_llvm_dialect_front.png}{矩阵乘法转换为 LLVM Dialect 的前半部分}{fig:mlir-llvm-dialect-front}

\labfigure{figures/mlir/mlir_llvm_dialect_address.png}{LLVM Dialect 中的地址计算、读写和浮点乘加}{fig:mlir-llvm-dialect-address}

\subsection{目标文件检查}

\texttt{matmul\_llvm.mlir} 属于 MLIR 的 LLVM Dialect。我们使用 \texttt{mlir-translate-14} 将其生成标准 LLVM IR：

\begin{lstlisting}[language=bash]
mlir-translate-14 --mlir-to-llvmir matmul_llvm.mlir -o matmul.ll
\end{lstlisting}

我们在 \texttt{matmul.ll} 中查看 \texttt{define}、\texttt{phi}、\texttt{icmp}、\texttt{br}、\texttt{getelementptr}、\texttt{load}、\texttt{fmul}、\texttt{fadd}、\texttt{store} 和 \texttt{ret}。\texttt{phi} 在基本块之间传递循环变量，乘法和加法将矩阵下标写成地址偏移。

\labfigure{figures/mlir/mlir_standard_llvm_front.png}{MLIR LLVM Dialect 生成的标准 LLVM IR 前半部分}{fig:mlir-standard-llvm-front}

\labfigure{figures/mlir/mlir_standard_llvm_loop.png}{标准 LLVM IR 中的循环结构与浮点乘加}{fig:mlir-standard-llvm-loop}

最后使用 Clang 将标准 LLVM IR 编译为目标文件：

\begin{lstlisting}[language=bash]
clang -c matmul.ll -o matmul.o
file matmul.o
nm -g matmul.o | grep ' matmul$'
\end{lstlisting}

我们得到 x86-64 ELF 可重定位文件 \texttt{matmul.o}，并在符号表中检查 \texttt{matmul}。我们使用 \texttt{clang -c -O0 -g} 编译仓库中的 \texttt{matmul.ll}，得到带 \texttt{matmul} 符号和 \texttt{.debug\_info} 节的目标文件。

\labfigure{figures/mlir/mlir_clang_object.png}{标准 LLVM IR 编译为 x86-64 可重定位目标文件}{fig:mlir-clang-object}

\subsection{本部分小结}

我们先对基础程序执行 Canonicalization，再以矩阵乘法为例依次记录 \texttt{linalg.matmul}、\texttt{scf.for}、基本块控制流、LLVM Dialect 和标准 LLVM IR。最后，我们生成 \texttt{matmul.o}，并在符号表中检查 \texttt{matmul}。
